% from ega4-21c-fr.tex
\oldpage[IV-4]{330}
tout ouvert $U$ de $X$, on a
$\mathscr M_{X/S}^*|U=\mathscr M_{U/S}^*$, donc
$\mathscr{D}iv_{X/S}|U=\mathscr{D}iv_{U/S}$, et le faisceau
$\mathscr{D}iv_{X/S}$ est donc égal au préfaisceau
$U\mapsto\operatorname{Div}(U/S)$.

Puisque $\mathscr M_{X/S}^*$ est un sous-faisceau de $\mathscr M_X^*$, les
définitions, notations et formules relatives aux diviseurs des sections de
$\mathscr M_X^*$ au-dessus de $X$ (21.1.3) s'appliquent sans changement aux sections
de $\mathscr M_{X/S}^*$ au-dessus de $X$.

\begin{env}[21.15.3]
\label{IV.21.15.3}
La structure de faisceau de groupes ordonnés sur $\mathscr{D}iv_X$ (21.1.6) induit
sur le sous-faisceau $\mathscr{D}iv_{X/S}$ une structure de faisceau de groupes
ordonnés, pour laquelle le faisceau de monoïdes des germes de sections positives est
$\mathscr{D}iv_X^+\cap\mathscr{D}iv_{X/S}$, que l'on note
$\mathscr{D}iv_{X/S}^+$. On a
$\Gamma(X,\mathscr{D}iv_{X/S}^+)=\operatorname{Div}^+(X)\cap
\operatorname{Div}(X/S)$; on note ce sous-monoïde de $\operatorname{Div}(X/S)$ par
$\operatorname{Div}^+(X/S)$, et il est formé des éléments $\geq0$ pour une structure
de groupe ordonné sur $\operatorname{Div}(X/S)$. Il résulte de (21.1.5.1) que
$\mathscr{D}iv_{X/S}^+$ est l'image dans $\mathscr{D}iv_{X/S}$ du sous-faisceau de
monoïdes
\[
\label{IV.21.15.3.1}
  \mathcal O_{X/S}^{\mathrm{reg}}=\mathcal O_X\cap\mathscr M_{X/S}^*.
  \tag{21.15.3.1}
\]
\end{env}

Pour tout ouvert $U$ de $X$, $\Gamma(U,\mathcal O_{X/S}^{\mathrm{reg}})$ est
l'ensemble des sections $t$ de $\mathcal O_X$ au-dessus de $U$ telles que $t$ soit
régulière et que $1/t$ appartienne à $\Gamma(U,\mathscr M_{X/S})$, ce qui signifie,
avec les notations de (20.6.1), que $t\in\mathscr T_{X/S}(U)$, de sorte que le faisceau
$\mathcal O_{X/S}^{\mathrm{reg}}$ n'est autre que le faisceau noté $\mathscr T_{X/S}$
dans (20.6.1). On peut donc écrire
\[
\label{IV.21.15.3.2}
  \mathscr{D}iv_{X/S}^+=\mathcal O_{X/S}^{\mathrm{reg}}/\mathcal O_X^*
  \tag{21.15.3.2}
\]
à un isomorphisme canonique près.

\begin{env}[21.15.3.3]
\label{IV.21.15.3.3}
Soit $D\in\operatorname{Div}^+(X/S)$ et considérons le sous-préschéma fermé $Y(D)$
de $X$ défini par l'Idéal $\mathscr J_X(D)$ de $\mathcal O_X$ (21.6.5); en vertu de
ce qui précède, pour tout $x\in Y(D)$, il y a un voisinage ouvert $U$ de $x$ dans $X$
et une section $t\in\mathscr T_{X/S}(U)$ telle que
$\mathscr J_X(D)|U=\mathcal O_U\cdot t$; puisque $x\in Y(D)$, l'image $t_{f(x)}$ de
$t$ dans $\Gamma(U\cap X_{f(x)},\mathcal O_{X_{f(x)}})$ appartient à l'idéal maximal
de $\mathcal O_{X_{f(x)},x}$ et en outre, par définition, pour tout $s\in S$, l'image
$t_s$ de $t$ dans $\Gamma(U\cap X_s,\mathcal O_{X_s})$ est une section
\emph{régulière}. On déduit donc de (11.3.8) et (19.2.4) que l'immersion canonique
$Y(D)\to X$ est \emph{transversalement régulière relativement à $S$ et de
codimension $1$} au point $x$. La réciproque étant immédiate, on voit qu'on peut
\emph{identifier canoniquement} les diviseurs positifs sur $X$ relatifs à $S$ aux
sous-préschémas fermés $Y$ de $X$ tels que l'injection canonique $Y\to X$ soit une
immersion transversalement régulière relativement à $S$ et de codimension $1$. Nous
ferons d'ordinaire cette identification.
\end{env}

\begin{proposition}[21.15.4]
\label{IV.21.15.4}
Soient $D$ un diviseur sur $X$, $\mathscr J_X(D)$ l'Idéal fractionnaire inversible
correspondant (21.2.5). Pour que $D\in\operatorname{Div}(X/S)$, il faut et il suffit
que, pour tout $x\in X$, on ait
$(\mathscr J_X(D)\cap\mathscr M_{X/S}^*)_x\ne0$ (ou, ce qui revient au même, que
l'on ait $\mathscr J_X(D)\subset\mathscr M_{X/S}$ et
$\mathscr J_X(D)^{-1}\subset\mathscr M_{X/S}$).
\end{proposition}

\begin{proof}
En effet, dire que $D\in\operatorname{Div}(X/S)$ signifie que pour tout $x\in X$, il
existe un voisinage ouvert $U$ de $x$ et une section
$f\in\Gamma(U,\mathscr M_{X/S}^*)$ tels que $D|U=\operatorname{div}_U(f)$; comme
$\mathscr J_X(D)|U$ est l'Idéal fractionnaire inversible $\mathcal O_Uf$, la
proposition en résulte aussitôt.
\end{proof}

